\section{Rubix: \textit{substrate} de ejecución unificado y Compliance-as-Code}

Apollo decide cómo debe cambiar el software; Rubix proporciona la base de ejecución sobre la que ocurren esos cambios. Aquí, \term{substrate} designa la capa de ejecución unificada que está por debajo de las aplicaciones y por encima de cada nube o hardware concreto. La documentación oficial actual de Palantir describe Rubix como una Kubernetes implementation hardened, autoscaling y highly available, que aloja AIP, Foundry, Apollo y los workload de los clientes. Su valor no está en ocultar todas las diferencias, sino en fijar las políticas de seguridad, de ciclo de vida y de disponibilidad que tienen que ser uniformes.

\subsection{Qué unifica Kubernetes y qué le sigue faltando}

\term{Kubernetes} es un sistema de orquestación de contenedores que se encarga de declarar y mantener recursos de workload, service, network y storage. Sus objetos nativos pueden describir replicas, volume e ingress, pero no conocen el operator access de un entorno de seguridad nacional, la vida útil de los nodos, las aprobaciones de cumplimiento normativo ni la semántica de un despliegue que abarca toda la fleet. Rubix construye sobre Kubernetes una API opinionated que codifica esa experiencia organizacional como comportamiento predeterminado.

\begin{figure}[H]
\centering
\includegraphics[width=0.94\textwidth]{images/05-rubix-layers.png}
\caption{Rubix se sitúa entre las aplicaciones y la infraestructura heterogénea; Apollo expresa versiones y políticas por encima de él.}
\figsource{Redibujado a partir de los subtítulos 15:00--20:00, `images/05-rubix-layers.png`.}
\end{figure}

\begin{knowledgebox}{Lectura de la figura: el límite de la abstracción es la «intención de la aplicación»}
El equipo de producto declara cuántas replica necesita, qué persistencia, qué endpoint expone y cómo se determina el quorum; después, el substrate lo traduce a provider-specific resource, certificate, network policy y lifecycle. Si la aplicación todavía tiene que conocer el ingress de cada nube, la PKI de cada centro de datos y cada checklist de cumplimiento normativo, la abstracción no ha liberado realmente la velocidad de experimentación.
\end{knowledgebox}

\begin{figure}[H]
\centering
\includegraphics[width=0.9\textwidth]{images/official-rubix-k8s.png}
\caption{La figura oficial de Palantir presenta Rubix como una implementación de Kubernetes hardened, autoscaling y highly available.}
\figsource{Documentación oficial de arquitectura de Palantir Rubix, consultada el 2026-08-11, `images/official-rubix-k8s.png`.}
\end{figure}

\begin{importantbox}{Términos clave: Multi-tenancy y aislamiento}
Multi-tenancy significa que varios clientes, equipos o workload comparten la infraestructura subyacente y, al mismo tiempo, conservan sus límites de identidad, datos, red y recursos. El aislamiento no se reduce al namespace; incluye también privilegio mínimo, network segmentation, gestión de secret, resource quota, auditoría y control de noisy-neighbor. Compartir aumenta la utilización, pero también amplía el blast radius de un error de configuración.
\end{importantbox}

\subsection{Del cumplimiento normativo manual a la ejecución por software}

Los controles regulatorios suelen tener objetivos razonables, por ejemplo restringir a los operator, registrar los cambios, cifrar las transmisiones y mantener una línea base de parches. Si dependen solo de documentos y de checklist manuales, los controles crecen linealmente con el número de entornos y tienden a producir drift. \term{Compliance-as-code} convierte los requisitos en policy, immutable configuration, automated evidence y gate que bloquean los despliegues que incumplen, de modo que «cumplir las reglas» sea un estado del sistema y no un conjunto de archivos reunidos a última hora para la auditoría trimestral.

\begin{figure}[H]
\centering
\includegraphics[width=0.94\textwidth]{images/official-rubix-security.png}
\caption{La figura oficial de seguridad de Rubix enumera controles geográficos, auditoría, zero-downtime upgrade, node cycling, encryption y service orchestration.}
\figsource{Documentación oficial de arquitectura de Palantir Rubix, consultada el 2026-08-11, `images/official-rubix-security.png`.}
\end{figure}

\begin{knowledgebox}{Lectura de la figura: las propiedades de seguridad deben traducirse en evidencia observable}
«Encryption at rest» exige poder demostrar la key y la storage policy; «network monitoring» exige registros consultables y alert; «geographic controls» exige restricciones en el scheduler y en la ubicación de los datos. El zero-downtime upgrade de la figura significa mantener el servicio durante una publicación; no es el ZeRO del entrenamiento distribuido, que es el Zero Redundancy Optimizer y ahorra memoria de cada GPU aplicando sharding al optimizer state, al gradient o a los parameter. La figura de seguridad es un catálogo de capacidades, no una prueba automática. Cada elemento debe vincularse con su configuración, su evidencia de ejecución, su manejo de anomalías y su owner.
\end{knowledgebox}

\subsection{Ephemeral node: eliminar el drift mediante reemplazo}

\term{Immutable image} significa que la image en ejecución no se modifica en el lugar; los parches entran al sistema construyendo una image nueva y reemplazando la instance. \term{Ephemeral node} es un nodo con vida útil limitada, del que se espera que sea destruido y reconstruido. Así se reduce el drift a largo plazo y se obliga a que los service tengan capacidad de failover. El node cycling también acorta el tiempo durante el cual un atacante puede mantener un foothold persistente, pero no sustituye a la credential rotation, la network policy ni la application security.

\begin{figure}[H]
\centering
\includegraphics[width=0.94\textwidth]{images/06-ephemeral-security.png}
\caption{La ephemerality convierte el patching en un ciclo controlado de replacement.}
\figsource{Redibujado a partir de los subtítulos y de la documentación actual de Rubix, `images/06-ephemeral-security.png`.}
\end{figure}

\begin{warningbox}{Diferencia de versiones: 72 horas en la exposición oral y 48 horas en la documentación actual}
En los subtítulos de la clase, Sankar dice que las container image se destruyen al azar en algún momento entre unas 40 y 72 horas y se reconstruyen a partir de una image nueva; la documentación oficial de Rubix consultada el 2026-08-11 dice que nodes cannot live longer than 48 hours. Puede que la política del producto se haya actualizado entre ambas fechas, o que la diferencia se deba a que se habla de image, host o node. Estas notas no fusionan ambas cifras en un solo número; al implementarlo, debe tomarse como referencia la policy vigente del despliegue de destino.
\end{warningbox}

\begin{table}[H]
\centering
\begin{tabular}{L{0.19\textwidth}L{0.32\textwidth}L{0.39\textwidth}}
\toprule
Mecanismo & Capacidad que se obtiene & Nuevos requisitos \\
\midrule
Immutable image & Línea base reproducible, menos drift en el lugar & Cadena de construcción y firma confiable \\
Bounded lifetime & Limita la persistencia, obliga a actualizar & Workload reprogramable, estado externalizado \\
Policy-driven drain & Menos inestabilidad por desalojos simultáneos & Visibilidad de topología, capacidad y prioridad de tareas \\
Blue/green rollout & Validación en paralelo, cambio zero-downtime & Capacidad duplicada y criterio de salud \\
Audit logging & Trazabilidad del comportamiento de operator y sistema & Integridad de los registros y control de acceso \\
\bottomrule
\end{tabular}
\caption{La automatización de la seguridad convierte las dependencias manuales implícitas en nuevos requisitos de software y de capacidad.}
\end{table}

\subsection{Resumen de la sección}

Rubix usa un substrate unificado de Kubernetes para encargarse de la seguridad, el aislamiento, el ciclo de vida y la uniformidad entre entornos; Apollo aprovecha ese substrate para ejecutar los cambios de versión y de configuración. Compliance-as-code y la ephemeral infrastructure reducen el toil manual, pero exigen que los workload admitan de verdad fallas, reconstrucción y estado externalizado. La abstracción no elimina la complejidad: la traslada a un equipo de plataforma que puede verificarla de forma centralizada.
